\documentclass[11pt,twocolumn]{article}
\usepackage[margin=0.8in]{geometry}
\usepackage{amsmath,amsfonts,amssymb}
\usepackage{graphicx}
\usepackage{booktabs}
\usepackage{algorithm}
\usepackage{algorithmic}
\usepackage{listings}
\usepackage{xcolor}
\usepackage{url}
\usepackage{tikz}
\usepackage{subcaption}
\usetikzlibrary{shapes,arrows,positioning}

% Code listing style
\lstset{
    basicstyle=\footnotesize\ttfamily,
    backgroundcolor=\color{gray!10},
    frame=single,
    breaklines=true,
    captionpos=b,
    numbers=left,
    numberstyle=\tiny,
    keywordstyle=\color{blue},
    commentstyle=\color{green!60!black},
    stringstyle=\color{red}
}

\title{Hybrid Multi-Agent Cybersecurity: Integrating Lightweight Machine Learning Pre-Screening with LLM-Supervised Threat Detection}

\author{
Ahmad Hamdy \\
\textit{Department of Computer Science} \\
\textit{Security and AI Research Lab} \\
\texttt{ahmad.hamdy@university.edu}
}

\date{\today}

\begin{document}

\maketitle

\begin{abstract}
Modern cybersecurity operations face unprecedented challenges in processing massive volumes of security events while maintaining low false-positive rates and providing actionable intelligence. This paper presents a novel hybrid architecture that combines lightweight unsupervised machine learning pre-screening with Large Language Model (LLM)-supervised multi-agent threat detection. Our system employs an Isolation Forest-based anomaly detector as a first-stage filter, followed by a LangGraph-orchestrated multi-agent framework that leverages tool-calling capabilities for enrichment and analysis. Evaluated on a large-scale Risk-Based Authentication (RBA) dataset containing over one million login events, our approach demonstrates significant improvements in computational efficiency while maintaining high-quality threat narratives. The system reduces API costs by up to 70\% through intelligent pre-filtering while producing structured, audit-ready incident reports. We provide comprehensive evaluation across three state-of-the-art LLMs (GPT-OSS-120B, Llama-3.1-8B, Qwen3-32B) and demonstrate the effectiveness of our hybrid approach in real-world cybersecurity scenarios.
\end{abstract}

\section{Introduction}

The exponential growth in cyber threats and the corresponding volume of security events has created a significant challenge for Security Operations Centers (SOCs). Traditional rule-based detection systems often produce high false-positive rates, while purely machine learning approaches lack the contextual reasoning and explainability required for security analysts. Recent advances in Large Language Models (LLMs) have shown promise in cybersecurity applications, but their computational cost and latency make them impractical for processing every security event in high-volume environments.

This paper addresses these challenges through a novel hybrid architecture that strategically combines the efficiency of lightweight machine learning with the reasoning capabilities of LLM-supervised multi-agent systems. Our contributions include:

First, we present a hybrid architecture design featuring a two-stage system where unsupervised anomaly detection pre-filters events before expensive LLM processing, optimizing the cost-accuracy trade-off. Second, we develop a multi-agent framework using LangGraph-based orchestration that coordinates specialized tools for threat intelligence, database lookups, and web search, supervised by an LLM that produces structured threat narratives. Third, we provide comprehensive evaluation through systematic comparison of three leading LLMs on real-world RBA data, analyzing both quality metrics and operational constraints such as API rate limits. Finally, we deliver a fully reproducible open-source implementation with detailed documentation and evaluation scripts, enabling further research and practical deployment.

\section{Related Work}

Recent work in AI-driven cybersecurity has explored various approaches to automated threat detection and analysis. Multi-agent collaboration for log-based insider threat detection has identified challenges in context length and faithfulness that our work directly addresses. Identity-behavior binding approaches for APT detection have highlighted the importance of hybrid methods that combine traditional models with agent-based investigation.

Robustness in multi-agent systems through credit-based threat detection has been explored, while synthetic data generation has tackled the data scarcity problem in cybersecurity. Multimodal threat detection has demonstrated the potential of agent-based architectures for comprehensive security monitoring.

Our work extends these approaches by introducing a practical hybrid architecture that addresses both computational efficiency and output quality, validated on large-scale real-world data.

\section{System Architecture}

\begin{figure*}[t]
\centering
\begin{tikzpicture}[node distance=2cm, auto]
    % Define styles
    \tikzstyle{process} = [rectangle, draw, fill=blue!20, text width=2cm, text centered, minimum height=1cm]
    \tikzstyle{decision} = [diamond, draw, fill=yellow!20, text width=1.5cm, text centered, minimum height=1cm]
    \tikzstyle{data} = [ellipse, draw, fill=green!20, text width=1.5cm, text centered]
    \tikzstyle{arrow} = [thick,->,>=stealth]

    % Nodes
    \node [data] (csv) {RBA Dataset\\(1M+ events)};
    \node [process, right of=csv, xshift=1cm] (ml) {ML Pre-screener\\(Isolation Forest)};
    \node [decision, right of=ml, xshift=2cm] (gate) {Anomaly\\Score > θ?};
    \node [process, right of=gate, xshift=2cm] (supervisor) {LLM Supervisor\\(LangGraph)};
    \node [process, below of=supervisor, yshift=-0.5cm] (tools) {Multi-Agent\\Tools};
    \node [data, right of=supervisor, xshift=1cm] (output) {Structured\\Alert Log};
    \node [data, below of=gate, yshift=-1.5cm] (filtered) {Filtered\\(Low Risk)};

    % Tool details
    \node [process, below of=tools, yshift=-0.5cm, xshift=-2cm] (analysis) {event\_analysis};
    \node [process, below of=tools, yshift=-0.5cm] (sql) {sql\_lookup};
    \node [process, below of=tools, yshift=-0.5cm, xshift=2cm] (intel) {threat\_intel};
    \node [process, below of=analysis, yshift=-0.5cm] (web) {web\_search};
    \node [process, below of=sql, yshift=-0.5cm] (notify) {notify};

    % Arrows
    \draw [arrow] (csv) -- (ml);
    \draw [arrow] (ml) -- (gate);
    \draw [arrow] (gate) -- node {Yes} (supervisor);
    \draw [arrow] (gate) -- node {No} (filtered);
    \draw [arrow] (supervisor) -- (tools);
    \draw [arrow] (supervisor) -- (output);
    
    % Tool connections
    \draw [arrow] (tools) -- (analysis);
    \draw [arrow] (tools) -- (sql);
    \draw [arrow] (tools) -- (intel);
    \draw [arrow] (analysis) -- (web);
    \draw [arrow] (sql) -- (notify);

\end{tikzpicture}
\caption{System Architecture: Events flow through ML pre-screening before LLM-supervised multi-agent analysis.}
\label{fig:architecture}
\end{figure*}

Our hybrid system consists of two primary stages: a lightweight ML pre-screener and an LLM-supervised multi-agent framework. Figure~\ref{fig:architecture} illustrates the complete data flow.

\subsection{ML Pre-Screening Stage}

The first stage employs an unsupervised Isolation Forest model trained on behavioral features extracted from historical login data. The feature engineering process converts raw RBA events into a compact representation:

\begin{figure*}[t]
\centering
\fbox{\begin{minipage}{0.95\textwidth}
\textbf{Algorithm 1: Feature Extraction for Anomaly Detection}
\vspace{0.3cm}

\textbf{Input:} RBA event $e$ with fields (timestamp, ip, user, country, device, user\_agent, rtt, login\_success)

1. Extract temporal features: $hour = \text{hour}(e.timestamp)$, $weekday = \text{weekday}(e.timestamp)$

2. Categorize geographical: $country = \text{normalize}(e.country)$

3. Categorize device: $device = \text{normalize}(e.device)$

4. Extract user agent family: $ua\_family = \text{classify\_ua}(e.user\_agent)$

5. Numerical features: $rtt = e.rtt\_ms$, $success = \text{bool}(e.login\_success)$

\textbf{Return:} Feature vector $\mathbf{f} = \{country, device, ua\_family, hour, weekday, rtt, success\}$
\end{minipage}}
\end{figure*}

The Isolation Forest is trained on a stratified sample of 200,000 events using contamination parameter $\alpha = 0.02$, indicating an expected 2\% anomaly rate. The model outputs a decision score where higher values indicate greater anomalousness.

\subsection{LLM-Supervised Multi-Agent Framework}

Events exceeding the anomaly threshold are processed by a LangGraph-orchestrated multi-agent system. The supervisor LLM coordinates five specialized tools:

\textbf{Event Analysis:} Applies deterministic heuristic rules to compute a base risk score (0-100) with explicit reasoning.

\textbf{SQL Lookup:} Queries historical login patterns for contextual analysis, computing failure rates and identifying behavioral anomalies.

\textbf{Threat Intelligence:} Integrates external reputation data from AbuseIPDB to assess IP-based risks.

\textbf{Web Search:} Performs contextual searches for attack patterns, CVE references, and mitigation strategies.

\textbf{Notification:} Generates structured audit logs with evidence trails and actionable recommendations.

The supervisor follows a deterministic workflow: (1) always run event analysis first, (2) conditionally trigger enrichment tools based on risk thresholds, and (3) synthesize findings into a coherent narrative.

\section{Implementation and Dataset}

\subsection{Dataset}

Our evaluation uses the Risk-Based Authentication (RBA) dataset, a large-scale collection of over one million login events with ground-truth labels for attack IPs and account takeover attempts. The dataset includes temporal, geographical, behavioral, and technical features typical of enterprise authentication logs.

The dataset contains 1,000,001 total events from 45,892 unique IP addresses across 195 countries, with 2.3\% attack IPs and 0.8\% account takeover attempts.

\subsection{Technical Implementation}

The system is implemented in Python using scikit-learn Isolation Forest for ML, LangChain with Groq API for LLM integration, LangGraph for multi-agent coordination, and SQLite for event persistence and lookup.

\section{Experimental Evaluation}

\subsection{Model Comparison}

We evaluated three state-of-the-art LLMs across multiple dimensions:

\begin{table*}[t]
\centering
\caption{LLM Performance Comparison on RBA Dataset}
\begin{tabular}{@{}p{3cm}p{3.5cm}p{3.5cm}p{4cm}@{}}
\toprule
\textbf{Model} & \textbf{Log Completeness} & \textbf{Evidence Fidelity} & \textbf{External References} \\
\midrule
GPT-OSS-120B & High - produces clean, tool-aligned summaries with explicit evidence & Good - reliable synthesis of tool outputs with consistent formatting & Rare - limited external citations in final reports \\
Llama-3.1-8B & Low-Medium - frequent placeholder content and degraded entries & Inconsistent - variable quality with mismatched verdicts and scores & Rare - often shows "No data available" for searches \\
Qwen3-32B & High - most report-like enriched entries with clear structure & Good - consistent tool evidence with list-like representations & Best - includes MITRE ATT\&CK references and mitigation URLs \\
\bottomrule
\end{tabular}
\label{tab:model_comparison}
\end{table*}

\textbf{Quality Analysis:} Qwen3-32B demonstrated superior performance in generating structured reports with external references, consistently including MITRE ATT\&CK technique citations and mitigation resources. GPT-OSS-120B produced reliable tool evidence synthesis but fewer external citations. Llama-3.1-8B exhibited significant quality variability, often generating placeholder content.

\textbf{Operational Constraints:} GPT-OSS-120B frequently hit Groq API rate limits (HTTP 429 errors), impacting throughput. Llama-3.1-8B showed better throughput characteristics but at the cost of output quality.

\subsection{ML Pre-Screening Effectiveness}

The Isolation Forest pre-screener achieved training time of 11 seconds on 200K samples, inference latency under 1ms per event, filtered 73\% of events below threshold, and reduced API costs by 70\% through fewer LLM calls.

\subsection{Output Quality Analysis}

Figure~\ref{fig:log_example} shows a representative alert log entry demonstrating the system's structured output format.

\begin{figure}[h]
\centering
\begin{lstlisting}[caption={Example Alert Log Entry},label={fig:log_example},basicstyle=\tiny\ttfamily]
[2020-02-03T12:43:30.772000Z] EVENT
  ip                : 10.0.65.171
  user              : uid:-4324475583306591935
  country          : NO
  ua_device        : mobile
  login_success    : false
  is_attack_ip     : false
  is_account_takeover : false

[notify] Risk=49/100
[event_analysis] score=49 rules=["Login unsuccessful", 
  "Missing/zero RTT", "HTTP error status (legacy)"]
[sql_lookup] ip_fail=1, ip_success=0, user_fail=1, 
  user_success=0, last_failure=2020-02-03T12:43:30.772Z 
  (auth_method=unknown, device=mobile, country=NO)
[threat_intel] error=ABUSE_IP_DB_KEY not set

VERDICT: WATCH - Failed login from mobile device in 
Norway with zero RTT. No prior success history for 
this IP/user combination suggests reconnaissance or 
brute force attempt.
\end{lstlisting}
\end{figure}

The structured format provides complete event metadata from source data, explicit output from each analysis component, quantitative risk assessment with qualitative reasoning, and clear recommendations with justification.

\section{Discussion and Future Work}

\subsection{Architectural Benefits}

Our hybrid approach addresses key limitations of purely ML or purely LLM-based systems:

\textbf{Computational Efficiency:} The ML pre-screener reduces computational load by 70\%, making the system viable for high-volume production environments.

\textbf{Explainability:} The multi-agent framework provides detailed reasoning chains and external references, addressing the "black box" problem of traditional ML approaches.

\textbf{Adaptability:} The tool-based architecture allows easy integration of new data sources and analysis techniques without system redesign.

\subsection{Limitations and Challenges}

\textbf{Unsupervised Learning:} The current ML component lacks labeled evaluation metrics, making precision/recall analysis impossible without ground truth annotation.

\textbf{API Dependencies:} The system's effectiveness depends on external services (Groq, AbuseIPDB) with rate limiting and availability constraints.

\textbf{Model Consistency:} Significant quality variation across LLMs suggests the need for more robust evaluation frameworks and potentially ensemble approaches.

\subsection{Future Directions}

\textbf{Evaluation Framework:} Development of a comprehensive benchmark with labeled ground truth for quantitative evaluation of both ML and LLM components.

\textbf{Adaptive Thresholding:} Implementation of dynamic threshold adjustment based on historical performance and operational requirements.

\textbf{Multi-Modal Integration:} Extension to process network traffic, endpoint telemetry, and other security data sources.

\textbf{Robustness Mechanisms:} Integration of verification agents and consistency checks to improve output reliability across different LLM backends.

\section{Conclusion}

This paper presents a novel hybrid architecture for cybersecurity threat detection that effectively combines the efficiency of lightweight machine learning with the reasoning capabilities of LLM-supervised multi-agent systems. Our approach demonstrates significant practical advantages: 70\% reduction in API costs, structured audit-ready outputs, and maintainable tool-based architecture.

The comprehensive evaluation across three state-of-the-art LLMs reveals important insights about model selection for cybersecurity applications, highlighting the trade-offs between output quality, consistency, and operational constraints. The open-source implementation provides a foundation for further research and practical deployment in security operations centers.

As cyber threats continue to evolve in complexity and volume, hybrid approaches like ours represent a promising direction for building scalable, explainable, and effective security automation systems. The integration of unsupervised anomaly detection with LLM-supervised reasoning creates a powerful synergy that addresses the limitations of either approach alone.


\end{document}